\documentclass[a4paper]{article}

\usepackage{graphicx}
\usepackage{amsmath,amssymb}
\usepackage{minted}
\usepackage{multirow}
\usepackage{float}
\usepackage{todonotes}
\usepackage{forest}
\usepackage{xstring}
\usepackage{xcolor}
\usepackage[most]{tcolorbox}
\usepackage{xparse}
\usepackage{natbib}

\usetikzlibrary{graphs,graphdrawing}
\usegdlibrary{layered}

% for external graphviz
\input{/home/donkarlo/Dropbox/repo/nd_language_ai_project/src/nd_language_ai/formal/computer/programming/tex/latex/code_clips/external_graphviz.tex}

% for tags
\input{/home/donkarlo/Dropbox/repo/nd_language_ai_project/src/nd_language_ai/formal/computer/programming/tex/latex/parts/control_sequence/kind/semtag/semtag.tex}

% for links
\input{/home/donkarlo/Dropbox/repo/nd_language_ai_project/src/nd_language_ai/formal/computer/programming/tex/latex/code_clips/seehere.tex}

\title{}
\date{}
\author{Mohammad Rahmani}

\begin{document}
    \maketitle
    
    
    \section{Die Präposition \textit{an}}
        
        \textbf{Grundidee:} \textit{an} bedeutet oft \textit{at}, \textit{on}, \textit{by} oder \textit{next to}. Die genaue Bedeutung hängt vom Kontext ab.
        
        \textbf{Basic idea:} \textit{an} often means \textit{at}, \textit{on}, \textit{by}, or \textit{next to}. The exact meaning depends on the context.
        
        \subsection{Neben oder bei etwas}
            
            \textbf{Deutsch:} \textit{an} kann bedeuten, dass jemand oder etwas neben, bei oder an einem bestimmten Ort ist.
            
            \textbf{English:} \textit{an} can mean that someone or something is next to, by, or at a certain place.
            
            \textbf{Beispiel:} Ich stehe an der Tür.
            
            \textbf{Translation:} I am standing at the door.
            
            \textbf{Beispiel:} Ich sitze am Tisch.
            
            \textbf{Translation:} I sit at the table.
        
        \subsection{Kontakt mit einer vertikalen Fläche}
            
            \textbf{Deutsch:} Bei einer Wand oder einer anderen vertikalen Fläche benutzt man oft \textit{an}.
            
            \textbf{English:} With a wall or another vertical surface, German often uses \textit{an}.
            
            \textbf{Beispiel:} Das Bild hängt an der Wand.
            
            \textbf{Translation:} The picture is hanging on the wall.
        
        \subsection{Zeitangaben}
            
            \textbf{Deutsch:} Bei manchen Zeitangaben benutzt man \textit{an}. Die Form \textit{am} bedeutet \textit{an dem}.
            
            \textbf{English:} With some expressions of time, German uses \textit{an}. The form \textit{am} means \textit{an dem}.
            
            \textbf{Beispiel:} am Morgen = an dem Morgen
            
            \textbf{Translation:} in the morning
            
            \textbf{Beispiel:} am Montag = an dem Montag
            
            \textbf{Translation:} on Monday
        
        \subsection{Bewegung zu einer Fläche oder Stelle}
            
            \textbf{Deutsch:} Mit Akkusativ kann \textit{an} eine Bewegung zu einer Fläche oder Stelle ausdrücken.
            
            \textbf{English:} With the accusative case, \textit{an} can express movement toward a surface or place.
            
            \textbf{Beispiel:} Ich hänge das Bild an die Wand.
            
            \textbf{Translation:} I hang the picture on the wall.
        
        \subsection{Zusammenfassung}
            
            \textbf{Deutsch:} \textit{an + Dativ} antwortet oft auf die Frage: Wo?
            
            \textbf{English:} \textit{an + dative} often answers the question: Where?
            
            \textbf{Deutsch:} \textit{an + Akkusativ} antwortet oft auf die Frage: Wohin?
            
            \textbf{English:} \textit{an + accusative} often answers the question: To where?

%            \bibliographystyle{plainnat}
%            \bibliography{/home/donkarlo/Dropbox/repo/nd_shared_research/bibliography/bibliography.bib}
\end{document}

